% ISIN und Datum sind Bezeichner, keine Kennzahlen:
% check-literals: skip
\begin{titlepage}
\centering
\vspace*{3cm}
{\Huge\bfseries GEA Group Aktiengesellschaft\par}
\vspace{0.8cm}
{\Large Anlageanalyse aus Prim\"arquellen\par}
\vspace{0.4cm}
{\large Xetra: G1A \textbar{} ISIN DE0006602006\par}
\vspace{2.5cm}
{\large Kurs \je{\Kurs}{EUR} am \KursDatum\par}
\vspace{0.3cm}
{\large B\"orsenwert \Mio{\Boersenwert}{EUR}\par}
\vspace{3cm}
{\large Stand: 20.~September 2026\par}
\vspace{1cm}
\begin{minipage}{0.86\linewidth}\small
\textbf{Die Frage, die dieser Bericht beantwortet:} Soll ein Anleger die Aktie der GEA
Group zum heutigen Kurs kaufen, und wenn nicht, zu welchem Kurs?

\medskip
\textbf{Farben.} Schwarz sind belegte Tatsachen und Rechnungen. \bk{Blau ist die
Einschätzung des Verfassers.} \vd{Rot ist das Votum.} Jede Zahl im Flie{\ss}text stammt
aus der Zahlenschicht \texttt{data/}; \GeprueftAnzahl{} gespeicherte Werte sind gegen die
gedruckte Seite gepr\"uft, die sie nennen.
\end{minipage}
\end{titlepage}
